\scanpage{152}
évidentes $\langle x_i,x'_i\rangle$ sur les $G_i\oplus\check G_i$, et $q_Q$ sur $Q$. \deleted{Alors} Soit $\overline D=\{\overline E_1,\ldots,\overline E_{\unsure{r|\nu}}\}$, avec $\overline E_i=\sum_{1\leq j\leq i}\unsure{R|P}_j\in\overline E$. \deleted{Considérons}

Considérons
\[
\left\{\begin{aligned}
\overline R_i&=\overline E_i/\overline E_{i-1}\\
\overline Q&=\overline E_r^\perp/\overline E_r
\end{aligned}\right.
\]
on a donc des isom canoniques
\[
\begin{aligned}
\overline R_i&\xrightarrow{\overset{\bar u_i}{\sim}}R_i\\
\overline Q&\mathrel{\overset{\unsure{\bar v|\bar u|\bar u_q}}{\simeq}}Q,
\end{aligned}
\]
le dernier étant compatible avec $\bar q,q$.

Ceci posé, on a

\uline{Cor. 5} L'ens $\Sigma$ est canoniquement isom à l'ens. des isom $(\overline E,\bar q,\overline D)\longrightarrow(\overset{\deleted{\unsure{-}}}{E},q,D)$ compatibles avec les $\bar u_i,\bar v$.

En termes schématiques, le schéma produit $\underline\Sigma=\prod_{1\leq\unsure{j|i}\leq r}\underline\Sigma_i\times\underline\Sigma_0$ est canoniquement isomorphe au schéma
\[
\underline{\mathrm{Isom}}_{\{\bar u_i\},\bar v}\big((\overline E,\bar q,\overline D);(E,q,D)\big)
\]
\margin{sans hyp. $S$ affine !}%
% slanted
Ce dernier est un torseur à gauche sous $\underline{\mathrm O}_D^+(q)$ (à droite sous $\underline{\mathrm O}_{\overline D}^{+}(\bar q)$)

Comme $\underline\Sigma$ est \unsure{évidemt|évidemment} lisse \inserted{$0$-}connexe sur $S$, on en conclut
% archive sheet number: 151
